% Front matter for the venue-fitted manuscript: abstract (AAS: 150-200 words),
% keywords, and the statements AAS requires. Not used by the 110-page paper/main.tex.

\begin{abstract}
Annuity and pension valuations price the \emph{spread} of a mortality forecast, not only its
centre. Nominal 95\% intervals have been produced since \citet{lee1992modeling}, lately by neural
extensions of Lee--Carter, but never audited against a structural break alongside a
quiet control that separates break-induced failure from standing miscalibration. We
pre-register that audit on the June 2026 Human Mortality Database vintage: ten forecasting
families crossed with seven uncertainty mechanisms, fifty admissible cells, one scoring path,
three regimes---an expanding-origin control (1990--2014), the COVID-19 break (train to 2019,
test 2020--2024) and a 1914--1922 placebo. The control carries the finding: most
miscalibration predates the pandemic. Native intervals from the Lee--Carter lineage already
cover 0.702, 0.682 and 0.731 against 0.950 in quiet decades, and the registered hypothesis of
universal under-coverage across the break is refuted---its two-sided restatement grows in only
27 of 50 cells, because the break mainly sorts arms by the side of nominal they began on.
Joint path coverage exceeds independence at each arm's own marginal rate in 50 of 50 control
and break cells. In money it is starker: the native Lee--Carter interval
covers 0.253 of realised annuity factors. Most were never calibrated; the pandemic mainly
re-orders them.
\end{abstract}

\noindent\textbf{Keywords:} mortality forecasting; longevity risk; calibration; prediction
intervals; proper scoring rules; conformal prediction; COVID-19; Lee--Carter; neural networks;
pre-registration.
